\documentclass[11pt]{article}
\usepackage[utf8]{inputenc}
\usepackage[margin=0.85in]{geometry}
\usepackage{setspace}
\setstretch{1.25}
\usepackage{hyperref}
\hypersetup{colorlinks=true, linkcolor=black, urlcolor=blue}

\begin{document}

\begin{center}
    {\large \textbf{RESEARCH HIGHLIGHTS}}\par
    {\footnotesize \textit{Document for Editorial Office --- Kept Separate from Double-Blind Anonymized Manuscript}}\par
    \vspace{0.4cm}
    \hrule height 1.2pt
\end{center}
\vspace{0.3cm}

\noindent \textbf{The Compute-Grid Collision: AI Load Shocks, Wholesale Capacity Rents, and the Regressive Ratepayer Tax} \\
\textit{Danil Popov --- Independent Researcher $\bullet$ Target Journal: Energy Policy} \\
\vspace{0.4cm}

\begin{itemize}
    \item \textbf{Inelastic AI Compute Demand:} Exorbitant GPU idle costs dominate electricity prices, rendering hyperscaler power demand hyper-inelastic and driving wholesale capacity clearing prices to administrative caps.
    \item \textbf{PJM Capacity Shock \& Pass-Through:} The July 2024 PJM auction triggered a +574\% procurement surge from \.2B to \.7B; capacity shocks pass into residential tariffs with a 0.38 elasticity (0.77 cents/kWh per \/MW-day).
    \item \textbf{FERC Co-Location Rejection:} Exploits the November 2024 FERC Order (Docket ER24-2172) rejecting behind-the-meter co-location, documenting an avoided \--\ annual transmission cost shift onto captive ratepayers.
    \item \textbf{Spatial Border Discontinuity:} Boundary RDD across contiguous county borders reveals an immediate +3.214 cents/kWh retail tariff jump and an 18.4\% manufacturing employment contraction ($-420$ jobs per county).
    \item \textbf{Regressive Ratepayer Incidence:} Bottom-quintile households incur a +2.38 percentage point energy burden expansion (+25.7\% relative increase) vs +0.21 pp (+14.8\%) for top-quintile earners (Regressivity Multiplier = 1.74x).
\end{itemize}

\end{document}

